\chapter{Faits empiriques des carnets d'ordres}\label{ch:mx:empirical-facts-of-order-books}

Au T2 2026, parmi les ordres inscrits dans les carnets des plus grandes actions américaines et non exécutés, 23~\% ont été annulés dans un
délai d'une milliseconde après leur arrivée, 54~\% dans un délai d'un dixième de seconde et 95~\% dans la minute (statistiques MIDAS de la SEC). Personne n'a conçu ce
nombre, et personne n'a conçu les autres que ce chapitre rassemble~: la profondeur affichée culmine derrière le meilleur prix, les nouveaux ordres sont placés
à des distances du meilleur prix qui suivent une loi de puissance, que la fourchette varie avec la \omterm{def:mx:empirical-facts-of-order-books:vol}{volatilité par transaction}. Ce sont les
faits stylisés du carnet d'ordres (le livre~7, chapitre~5, pour ceux des rendements)~: des régularités qui se vérifient pour les actions, les plateformes et les années, que tout modèle
de carnet doit reproduire, et que le marché simulé utilisé dans ce livre ne les reproduit qu'en partie. Le chapitre mesure les deux et dit quels faits
le simulateur ne reproduit pas.

\section{Profils de profondeur}

Le carnet moyen d'une action liquide n'est pas le plus profond au meilleur prix. Bouchaud, Mézard et Potters (2002) ont observé, sur trois actions cotées à la Bourse de Paris, un
profil moyen en bosse~: la profondeur croît pendant quelques ticks derrière le meilleur prix puis décroît, forme que reproduit un \omterm{def:mx:the-limit-order-book:zi}{modèle à intelligence nulle}.
L'heure simulée du chapitre~1 présente la même forme (la \cref{fig:mx:the-limit-order-book:fill}, à droite)~: 1\,366 actions au meilleur prix en moyenne, 1\,508 à un
tick derrière. La raison apparaît dans les flux. Sur la journée simulée (6.5 heures, 715\,054 messages), les annulations par action affichée et par seconde
sont de 0.090 au meilleur prix, 0.057 à un tick derrière et 0.044 à deux ticks derrière~: le meilleur niveau est celui où les ordres deviennent obsolètes en premier et sont exécutés en premier, si bien que
il porte moins de liquidité que les niveaux qu'un \omterm{def:mx:the-limit-order-book:orders}{ordre au marché} atteint rarement.

\section{Où vont les nouveaux ordres}

\begin{definition}[Prix limite relatif]\label{def:mx:empirical-facts-of-order-books:relative}
Le \emph{prix limite relatif}\index{prix limite relatif} d'un nouvel \omterm{def:mx:the-limit-order-book:orders}{ordre à cours limité} est sa distance, en ticks, au meilleur prix de son côté à son
arrivée~: positive derrière le meilleur prix, nulle à celui-ci, négative lorsqu'il l'améliore.
\end{definition}

Zovko et Farmer (2002) ont rassemblé sept millions d'ordres de la London Stock Exchange portant sur 50 actions et ont constaté que la distribution cumulée des \omterm{def:mx:empirical-facts-of-order-books:relative}{prix limites
relatifs} décroît à peu près comme une loi de puissance d'exposant voisin de 1.5, de quelques ticks à environ deux mille~: la plupart des ordres arrivent près du meilleur prix, mais quelques-uns
sont placés très loin, et leur part décroît lentement. Bouchaud, Mézard et Potters ont aussi trouvé une loi de puissance. Le simulateur n'en a pas~: le
\texttt{firm.tape} du livre~7 place des \omterm{def:mx:the-limit-order-book:orders}{ordres à cours limité} sur ses dix meilleurs niveaux, à des taux décroissant géométriquement. Dans la journée simulée, 0.5~\% des nouveaux ordres améliorent le
meilleur prix, 18.6~\% le rejoignent, 40.1~\% se placent à quatre ticks ou plus derrière, et aucun ne dépasse neuf. Une stratégie dont l'edge repose sur une liquidité profonde (un
grand métaordre qui parcourt le carnet, chapitre~11) trouverait le carnet simulé trop mince loin du meilleur prix.

\section{Fourchettes et volatilité de chaque transaction}

\begin{definition}[Volatilité par transaction]\label{def:mx:empirical-facts-of-order-books:vol}
La \emph{volatilité par transaction}\index{volatilité par transaction} $\sigma_1$ d'un instrument est l'écart-type de la variation du prix milieu entre
une transaction et la suivante.
\end{definition}

Wyart, Bouchaud, Kockelkoren, Potters et Vettorazzo (2008) ont montré que si ni le market making ni le market taking ne rapporte plus qu'un profit marginal,
la fourchette doit être proportionnelle à l'impact instantané d'une transaction~; ils ont trouvé que la fourchette et la \omterm{def:mx:empirical-facts-of-order-books:vol}{volatilité par transaction} sont fortement corrélées
d'une action à l'autre, avec un $R^2$ supérieur à 0.9, et y ont vu la preuve que la sélection adverse fixe la fourchette et que l'essentiel de la volatilité provient de l'impact
des transactions (les chapitres 4, 5 et 11 reprennent tour à tour les trois moitiés de cette phrase). La \Cref{fig:mx:empirical-facts-of-order-books:spread} teste la relation
dans deux marchés simulés. Dans le marché à intelligence nulle du chapitre~1, dont la fourchette peut varier librement, huit configurations des taux s'alignent sur
$s \approx 0.96 + 1.64\,\sigma_1$ et $R^2 = 0.989$. Dans dix configurations de \texttt{firm.tape}, la relation disparaît presque ($R^2=0.48$, pente 0.19)~:
l'action simulée est une action à grand tick dont la fourchette vaut un tick 96.6~\% du temps, et la fourchette n'a aucune marge pour réagir. Le chapitre~6 nomme les deux
régimes~; la relation est un fait propre aux marchés à petit tick.

\begin{omfigure}
\begin{tikzpicture}
\begin{axis}[width=9cm,height=5.8cm,xlabel={\small volatilité par transaction $\sigma_1$ (ticks)},ylabel={\small fourchette moyenne (ticks)},
  tick label style={font=\footnotesize},xmin=0,xmax=5.5,ymin=0,ymax=10,grid=major,grid style={gray!25},
  legend columns=3,legend style={font=\footnotesize,at={(0.5,-0.3)},anchor=north,draw=none,fill=none,
  /tikz/every even column/.append style={column sep=8pt}},table/col sep=comma]
\addplot[only marks,mark=o,omDef,thick] table[x=vol,y=spread,restrict expr to domain={\thisrow{kind}}{0:0}]{figdata/microstructure/03-empirical-facts-of-order-books/spreadvol.csv};
\addplot[only marks,mark=square*,omThm] table[x=vol,y=spread,restrict expr to domain={\thisrow{kind}}{1:1}]{figdata/microstructure/03-empirical-facts-of-order-books/spreadvol.csv};
\addplot[omDef,thick,dashed] table[x=vol,y=spread]{figdata/microstructure/03-empirical-facts-of-order-books/spreadvol_fit.csv};
\legend{marché à intelligence nulle,\texttt{firm.tape},ajustement $0.96+1.64\,\sigma_1$}
\end{axis}
\end{tikzpicture}

\omcaption{Fourchette moyenne en fonction de la \omterm{def:mx:empirical-facts-of-order-books:vol}{volatilité par transaction}~: huit configurations du marché à intelligence nulle du chapitre~1 (100\,000 événements chacune) et dix de
\texttt{firm.tape} (20 minutes simulées chacune, intensité des traders informés et taux de sauts du prix efficient variés). Les fourchettes du marché à intelligence nulle suivent la
droite~; celles de \texttt{firm.tape} restent proches d'un tick. Données~: \texttt{mx\_facts.spread\_vs\_vol\_zi}, \texttt{mx\_facts.spread\_vs\_vol}.}\label{fig:mx:empirical-facts-of-order-books:spread}
\end{omfigure}

\section{Formes intrajournalières}

Le volume, les messages et la volatilité suivent l'horloge (le livre~7, chapitre~5)~: élevés à l'ouverture, plus faibles à midi, de nouveau élevés à l'approche de la clôture, tandis que
les fourchettes des carnets réels sont les plus larges dans les premières minutes et se resserrent au fil de la matinée. \texttt{firm.tape} n'en reproduit que la première moitié~: ses taux sont
multipliés par un profil en U. Une journée simulée est trop bruitée pour le montrer, car un niveau d'activité aléatoire persistant domine chaque heure~; en moyenne sur
huit séances compressées, le volume passe de 92\,200 actions dans le premier treizième de la séance à 24\,700 au plus bas, puis remonte à 55\,600 dans le
dernier (voir la \cref{fig:mx:empirical-facts-of-order-books:intraday}). Sa fourchette reste à un tick à chaque heure~: le simulateur n'a pas de fourchette d'ouverture.

\begin{omfigure}
\begin{tikzpicture}
\begin{axis}[width=8.8cm,height=5cm,xlabel={\small treizième de la séance},ylabel={\small milliers},tick label style={font=\footnotesize},
  xtick={1,2,3,4,5,6,7,8,9,10,11,12,13},ymin=0,ymax=100,grid=major,grid style={gray!25},ybar,bar width=6pt,
  legend columns=2,legend style={font=\footnotesize,at={(0.5,-0.32)},anchor=north,draw=none,fill=none,
  /tikz/every even column/.append style={column sep=8pt}},table/col sep=comma]
\addplot[fill=omDef!60,draw=omDef] table[x=bin,y=volume]{figdata/microstructure/03-empirical-facts-of-order-books/intraday.csv};
\addplot[fill=omThm!60,draw=omThm] table[x=bin,y=messages]{figdata/microstructure/03-empirical-facts-of-order-books/intraday.csv};
\legend{actions échangées,messages}
\end{axis}
\end{tikzpicture}

\omcaption{Profil intrajournalier simulé~: actions échangées et messages dans chaque treizième de séance, moyennés sur huit séances de \texttt{firm.tape}
avec un profil d'activité en U. Données~: \texttt{mx\_facts.intraday\_mean}.}\label{fig:mx:empirical-facts-of-order-books:intraday}
\end{omfigure}

\section{Durées de vie des ordres, ordres fugaces et résilience du carnet}

\begin{definition}[Durée de vie de l'ordre, ordre fugace]\label{def:mx:empirical-facts-of-order-books:life}
La \emph{durée de vie de l'ordre}\index{durée de vie de l'ordre} est le temps écoulé entre son arrivée et son annulation totale ou son exécution complète. Un \emph{ordre
fugace}\index{ordre fugace} est un ordre intégralement annulé dans un bref délai après son arrivée (deux secondes constituent un seuil fréquent).
\end{definition}

Hasbrouck et Saar (2009) ont étudié de tels ordres sur des actions cotées au Nasdaq, sur l'Island ECN, un \omterm{def:mx:the-limit-order-book:lob}{carnet d'ordres à cours limité} électronique, et ont soutenu qu'ils se comportent davantage comme des demandes de liquidité que comme des offres
de celle-ci~: les traders affichent puis retirent leurs ordres pour chercher des contreparties cachées ou pour suivre un prix en mouvement. Le système MIDAS de la SEC publie les durées de vie de chaque
action américaine~: le délai entre l'arrivée d'un ordre et son annulation ou son exécution, l'horloge repartant après chaque événement partiel, estimé comme une table de survie
dans laquelle les exécutions censurent les annulations. La \Cref{fig:mx:empirical-facts-of-order-books:life} compare les courbes d'annulation à celles du simulateur. Pour
les grandes actions, la moitié des annulations survient dans un dixième de seconde, et dans les petites actions en environ une seconde~; dans \texttt{firm.tape}, 6.8~\% surviennent en
une seconde et 40~\% en dix secondes. Les ordres du simulateur vivent cent fois plus longtemps que les ordres réels~: ses participants sont lents, et les chapitres qui dépendent de
la course entre annulations et exécutions (latence, chapitre~8~; positions dans la file, chapitre~6) doivent préciser sur quel marché ils sont mesurés.

\begin{omfigure}
\begin{tikzpicture}
\begin{axis}[width=9cm,height=5.4cm,xlabel={\small temps depuis l'arrivée ou le dernier événement partiel},ylabel={\small part annulée},
  tick label style={font=\footnotesize},xtick={0,1,2,3,4,5},xticklabels={1 ms,10 ms,100 ms,1 s,10 s,1 min},ymin=0,ymax=1,
  grid=major,grid style={gray!25},legend columns=2,legend style={font=\footnotesize,at={(0.5,-0.3)},anchor=north,draw=none,fill=none,
  /tikz/every even column/.append style={column sep=8pt}},table/col sep=comma]
\addplot[omDef,very thick,mark=*] table[x=k,y=large]{figdata/microstructure/03-empirical-facts-of-order-books/lifetimes.csv};
\addplot[omProp,thick,mark=triangle*] table[x=k,y=mid]{figdata/microstructure/03-empirical-facts-of-order-books/lifetimes.csv};
\addplot[omThm,thick,mark=square*] table[x=k,y=small]{figdata/microstructure/03-empirical-facts-of-order-books/lifetimes.csv};
\addplot[black,thick,dashed,mark=o] table[x=k,y=sim]{figdata/microstructure/03-empirical-facts-of-order-books/lifetimes.csv};
\legend{grandes actions américaines,actions américaines de moyenne capitalisation,petites actions américaines,\texttt{firm.tape} (journée simulée)}
\end{axis}
\end{tikzpicture}

\omcaption{La part des ordres annulés dans un délai donné, les exécutions étant traitées comme des censures (estimations par table de survie)~: actions américaines par groupe de capitalisation en
2026 T2 (SEC MIDAS, Hazards and Survivors) et journée simulée. Données~: \texttt{data/microstructure/midas\_lifetimes.csv}, \texttt{mx\_facts.lifetable}.}\label{fig:mx:empirical-facts-of-order-books:life}
\end{omfigure}

\begin{definition}[Résilience du carnet]\label{def:mx:empirical-facts-of-order-books:resilience}
La \emph{résilience du carnet}\index{résilience du carnet} d'un \omterm{def:mx:the-limit-order-book:lob}{carnet d'ordres à cours limité} est la vitesse à laquelle son écart et sa profondeur reviennent à leur état habituel après qu'une
transaction ou une annulation les a épuisés.
\end{definition}

Large (2007) l'a mesurée sur un carnet électronique comme la réponse des nouveaux ordres aux chocs de liquidité. Dans la journée simulée, 750 exécutions ont vidé le meilleur
niveau et élargi la fourchette à deux ticks ou plus~; le temps médian avant que la fourchette ne revienne à un tick a été de 0.19 seconde, et 92.4~\% se sont rétablis
en moins d'une seconde. La résilience est ce qui distingue une concession de prix temporaire d'une concession permanente, et les noyaux d'impact du chapitre~12 sont la résilience mesurée en
prix.

Les faits varient aussi avec la structure de marché, et les statistiques MIDAS montrent l'ampleur de ce changement en une décennie.

\begin{dated}{2026-06}{Flux d'ordres sur les actions américaines dans les statistiques MIDAS}\label{dat:mx:empirical-facts-of-order-books:midas}
Moyennes quotidiennes sur l'ensemble des bourses américaines pour les actions (SEC MIDAS, Summary Metrics by Exchange)~: le ratio annulations sur transactions valait 19.8 en 2012 et 16.6 au
premier semestre 2026~; le taux caché (la part des transactions exécutées contre des \omterm{def:mx:order-types-and-their-uses:hidden}{ordres cachés}) est passé de 10.5~\% à 32.2~\%~; le taux de lots irréguliers (la part des transactions
inférieures à un lot standard) de 19.2~\% à 68.4~\%. Au T2 2026, parmi les ordres portant sur les plus grandes actions, 54~\% des annulations sont intervenues dans les 100 millisecondes
suivant l'arrivée et 95~\% dans la minute.
\end{dated}

\begin{omfigure}
\begin{tikzpicture}
\begin{axis}[name=rates,width=7cm,height=5cm,title={\small part des transactions (\%)},tick label style={font=\footnotesize},
  xtick={2012,2014,2016,2018,2020,2022,2024,2026},x tick label style={/pgf/number format/1000 sep={},rotate=45,anchor=east},
  ymin=0,ymax=75,grid=major,grid style={gray!25},legend columns=2,legend style={font=\footnotesize,at={(1.1,-0.35)},anchor=north,draw=none,
  fill=none,/tikz/every even column/.append style={column sep=8pt}},table/col sep=comma]
\addplot[omDef,very thick,mark=*] table[x=year,y=oddlot]{figdata/microstructure/03-empirical-facts-of-order-books/yearly.csv};
\addplot[omThm,thick,mark=square*] table[x=year,y=hidden]{figdata/microstructure/03-empirical-facts-of-order-books/yearly.csv};
\legend{taux de lots irréguliers,taux caché}
\end{axis}
\begin{axis}[at={(rates.east)},anchor=west,xshift=1.2cm,width=7cm,height=5cm,title={\small ratio annulations sur transactions},tick label style={font=\footnotesize},
  xtick={2012,2014,2016,2018,2020,2022,2024,2026},x tick label style={/pgf/number format/1000 sep={},rotate=45,anchor=east},
  ymin=0,ymax=25,grid=major,grid style={gray!25},table/col sep=comma]
\addplot[black,thick,mark=o] table[x=year,y=ctt]{figdata/microstructure/03-empirical-facts-of-order-books/yearly.csv};
\end{axis}
\end{tikzpicture}

\omcaption{Actions américaines, moyennes annuelles des statistiques quotidiennes sur l'ensemble des bourses, de 2012 à juin 2026~: la part des transactions en lots irréguliers et contre des \omterm{def:mx:order-types-and-their-uses:hidden}{ordres cachés}
(à gauche), et les annulations par transaction (à droite). Données~: SEC MIDAS, Summary Metrics by Exchange (\texttt{data/microstructure/midas\_yearly.csv}).}\label{fig:mx:empirical-facts-of-order-books:yearly}
\end{omfigure}

\section{Le simulateur est-il un marché~?}

Le tableau place les faits du chapitre à côté des chiffres du simulateur. Les 15.5 annulations par transaction de la journée simulée coïncident avec le ratio MIDAS des
dernières années~; la plupart des autres ne coïncident pas.

\begin{center}
\small
\begin{tabular}{@{}p{3.6cm}p{4.6cm}p{3.8cm}l@{}}
\toprule
fait & carnets réels & \texttt{firm.tape} & vérifié~? \\
\midrule
profil de profondeur & en bosse (Bourse de Paris) & pic à un tick derrière le meilleur & oui \\
\omterm{def:mx:empirical-facts-of-order-books:relative}{prix limites relatifs} & loi de puissance, exposant $\approx1.5$, jusqu'à 2\,000 ticks (Londres) & aucun au-delà de neuf ticks & non \\
écart et \omterm{def:mx:empirical-facts-of-order-books:vol}{volatilité par transaction} & $R^2>0.9$ selon les actions & $R^2=0.48$ (tick large~; le marché à intelligence nulle~: 0.989) & non \\
volume intrajournalier & en U & en U en moyenne (imposé) & oui \\
écart à l'ouverture & le plus large en début de journée & un tick toute la journée & non \\
vitesse d'annulation & 54~\% en moins de 0.1 s (grandes actions américaines) & 0.8~\% en moins de 0.1 s & non \\
ratio annulations sur transactions & 16.6 (actions américaines, 2026) & 15.5 & oui \\
liquidité cachée & 32~\% des transactions (2026) & aucune & non \\
résilience après une transaction qui vide le carnet & rapide (moins d'une seconde pour les actions liquides) & médiane 0.19 s & oui \\
\bottomrule
\end{tabular}
\end{center}

Un simulateur destiné à mesurer les positions dans la file, la latence ou la liquidité cachée doit d'abord corriger ses lignes en échec~; le chapitre~27 construit un marché à agents
calibré sur ces faits avec \texttt{firm.lobstats}.

\section{Tutoriel~: mesurer les faits}\label{tut:mx:empirical-facts-of-order-books:tut}

\begin{tutorial}
\textbf{Objectif.} Mesurer les faits stylisés sur la journée simulée et les placer à côté des statistiques de la SEC. \textbf{État final~:} le tableau et les
\cref{fig:mx:empirical-facts-of-order-books:spread,fig:mx:empirical-facts-of-order-books:life}.
\begin{tutsteps}
\item \textbf{Les données publiques.} \texttt{mx\_fetch\_midas.py} télécharge les fichiers MIDAS de séquences et de bourses pour le T2 2026 et écrit les deux tables dérivées dans
\texttt{data/microstructure/} (les copies versionnées suffisent pour la suite).
\item \textbf{Durées de vie.} La table de survie avec les exécutions comme censures, telle que MIDAS la calcule~:
\omcode{code/microstructure/03-empirical-facts-of-order-books/python/mx_facts.py}{121}{141}{Estimation par table de survie du délai avant annulation.}\label{lst:mx:empirical-facts-of-order-books:life}
\item \textbf{Flux et formes.} \texttt{relative\_prices}, \texttt{cancel\_rates}, \texttt{spread\_profile}, \texttt{intraday\_mean}, \texttt{resilience}
et \texttt{fleeting\_share} sur \texttt{day()}.
\item \textbf{La relation.} \texttt{spread\_vs\_vol\_zi()} et \texttt{spread\_vs\_vol()}~; tracer avec \texttt{fig\_facts.py}.
\end{tutsteps}
\textbf{Pour aller plus loin.} Donner à \texttt{firm.tape} un tick d'un demi-cent (diviser tous les prix par deux) et relancer la relation écart--volatilité~; remplacer son placement géométrique
par une loi de puissance et mesurer à nouveau le profil de profondeur.
\end{tutorial}

\section{Construction~: la boîte à outils des faits stylisés}\label{bld:mx:empirical-facts-of-order-books:lobstats}

\begin{build}
\bfield{Rôle} Un rapport unique sur les faits stylisés d'un carnet d'ordres, exécuté sur n'importe quel flux market-by-order~: celui du simulateur, un flux enregistré, un flux généré
(les données synthétiques du livre~12)~; la cible à laquelle le chapitre~27 se calibre.
\bfield{Interface} \texttt{firm\_lobstats.report(tape, levels, points)} renvoie le profil de profondeur, la distribution des prix relatifs,
les taux d'annulation selon la distance, la distribution des écarts, la \omterm{def:mx:empirical-facts-of-order-books:vol}{volatilité par transaction}, les courbes de survie des durées (annulations et exécutions), la part des \omterm{def:mx:empirical-facts-of-order-books:life}{ordres fugaces},
la résilience et le ratio annulations sur transactions~; \texttt{kaplan\_meier(durations, observed, points)}~; \texttt{compare(report, targets)} qui vérifie chaque fait
par rapport à un intervalle.
\bfield{Règles} Les entrées sont au format de messages de \texttt{firm.tape} (la fonction \texttt{res.tape()} du simulateur le produit)~; les durées de vie repartent à zéro après les événements partiels et
les exécutions censurent les annulations, comme dans MIDAS~; moyennes pondérées par le temps pour l'écart et la profondeur.
\bfield{Tests d'acceptation} \texttt{code/firm/lobstats/tests/}~: l'estimateur produit-limite à la main~; la profondeur en bosse d'une session simulée, un écart d'un tick, des tables de survie
monotones et une résilience rapide~; \texttt{compare} signalant les annulations lentes du simulateur.
\bfield{Pour aller plus loin} Un ajustement de queue en loi de puissance avec l'estimateur de Hill (livre~4, chapitre~15) et son intervalle de confiance~; le même rapport sur une
journée enregistrée issue de \texttt{make\_recorded\_day.py}.
\end{build}

\begin{omsources}
\item J.-P.~Bouchaud, M.~Mézard et M.~Potters, ``Statistical properties of stock order books: empirical results and models'', \emph{Quantitative
Finance} 2, 2002.
\item I.~Zovko et J.~D.~Farmer, ``The power of patience: a behavioural regularity in limit-order placement'', \emph{Quantitative Finance} 2, 2002.
\item M.~Wyart, J.-P.~Bouchaud, J.~Kockelkoren, M.~Potters et M.~Vettorazzo, ``Relation between bid-ask spread, impact and volatility in order-driven
markets'', \emph{Quantitative Finance} 8(1), 2008.
\item J.~Hasbrouck et G.~Saar, ``Technology and liquidity provision: the blurring of traditional definitions'', \emph{Journal of Financial Markets}
12(2), 2009.
\item A.~J.~Large, ``Measuring the resiliency of an electronic limit order book'', \emph{Journal of Financial Markets} 10(1), 2007.
\item B.~Biais, P.~Hillion et C.~Spatt, ``An empirical analysis of the limit order book and the order flow in the Paris Bourse'', \emph{Journal of
Finance} 50(5), 1995.
\item SEC, MIDAS market structure data: Hazards and Survivors by Time Period; Summary Metrics by Exchange; quote-life methodology.
\end{omsources}

\section{Exercices}

\begin{exercise}[$\star$]\label{exo:mx:empirical-facts-of-order-books:1}
Un ordre d'achat à cours limité arrive à 99.97 alors que le meilleur cours acheteur est de 100.00 et le meilleur cours vendeur de 100.01, avec un tick d'un centime. Calculer son \omterm{def:mx:empirical-facts-of-order-books:relative}{prix limite relatif}, puis celui d'un achat
à 100.00.
\end{exercise}

\begin{exercise}[$\star$]\label{exo:mx:empirical-facts-of-order-books:2}
Sur un marché à petit tick, la \omterm{def:mx:empirical-facts-of-order-books:vol}{volatilité par transaction} vaut 0.9 tick. Déterminer la fourchette que prédit le \omterm{def:mx:the-limit-order-book:zi}{modèle à intelligence nulle} ajusté dans le chapitre.
\end{exercise}

\begin{exercise}[$\star$]\label{exo:mx:empirical-facts-of-order-books:3}
Expliquer pourquoi le taux caché de la \cref{fig:mx:empirical-facts-of-order-books:yearly} est une part des transactions et non une part des ordres, et ce que
mesurerait à la place la part des ordres.
\end{exercise}

\begin{exercise}[$\star\star$]\label{exo:mx:empirical-facts-of-order-books:4}
Un trader observe 1\,000 annulations sur une grande action. D'après la table de survie du T2 2026, estimer combien d'entre elles surviennent dans les 100 millisecondes
suivant l'arrivée des ordres, puis dire ce que change la lecture «~exécutions traitées comme censure~».
\end{exercise}

\begin{exercise}[$\star\star$]\label{exo:mx:empirical-facts-of-order-books:5}
Le taux de lots irréguliers est passé de 19~\% des transactions en 2012 à 68~\% en 2026. Donner deux raisons structurelles et une conséquence pour un gestionnaire de flux qui ignore les lots
irréguliers.
\end{exercise}

\begin{exercise}[$\star\star$]\label{exo:mx:empirical-facts-of-order-books:6}
Expliquer pourquoi une action à grand tick brise la relation fourchette--volatilité, et dire ce qui remplace la fourchette comme variable d'ajustement.
\end{exercise}

\begin{exercise}[$\star\star\star$]\label{exo:mx:empirical-facts-of-order-books:7}
\emph{Programmation.} Calculer la table de survie du temps jusqu'à l'\emph{exécution} sur la journée simulée (\texttt{lifetable} avec les exécutions comme événement, les annulations
en censure) à 1, 10 et 60 secondes, puis la comparer aux chiffres de MIDAS pour les grandes actions~: 3.4~\%, 8.7~\% et 15.0~\%. Que dit la comparaison~?
\end{exercise}

\begin{exercise}[$\star\star\star$]\label{exo:mx:empirical-facts-of-order-books:8}
\emph{Trouver la faille.} «~Notre stratégie passive exécute 12~\% de ses ordres dans le simulateur et ne perd rien face aux traders plus rapides, elle fera donc la même chose en conditions réelles.~»
\end{exercise}

\section{Problème~: le simulateur est-il un marché~?}

\begin{problem}[{Problème du week-end --- le simulateur est-il un marché~?}]\label{pb:mx:empirical-facts-of-order-books:1}
Avant de se fier à \texttt{firm.tape} pour étudier l'exécution passive, un chercheur le confronte aux faits publiés et aux statistiques de la SEC.

\medskip\noindent\textbf{Partie I --- Le carnet.}
\begin{enumerate}
\item Où le profil de profondeur simulé atteint-il son maximum, et que montrent les carnets réels~?
\item Quels sont les taux d'annulation simulés par action affichée au meilleur cours et à deux ticks derrière, et pourquoi diffèrent-ils~?
\item Quelle part des nouveaux ordres améliore le meilleur cours, s'y aligne, ou se place quatre ticks ou plus derrière~?
\item À quoi ressemblent les \omterm{def:mx:empirical-facts-of-order-books:relative}{prix limites relatifs} réels, et que fait le simulateur à la place~?
\end{enumerate}

\medskip\noindent\textbf{Partie II --- Fourchettes et horloges.}
\begin{enumerate}[resume]
\item À quelle fréquence la fourchette simulée vaut-elle un tick, et quelle est sa moyenne~?
\item Quelle relation entre fourchette et \omterm{def:mx:empirical-facts-of-order-books:vol}{volatilité par transaction} donne le marché à intelligence nulle, avec quel $R^2$~?
\item Que donne \texttt{firm.tape}, et pourquoi~?
\item Quel profil de volume intrajournalier le simulateur produit-il, et combien de séances ont-elles été nécessaires pour le voir~?
\end{enumerate}

\medskip\noindent\textbf{Partie III --- Durées de vie.}
\begin{enumerate}[resume]
\item Comment MIDAS calcule-t-il ses courbes de durée de vie~?
\item Quelle part des annulations sur les grandes actions américaines, et sur les petites, survient dans les 100 millisecondes et dans la seconde~?
\item Quelles sont les parts du simulateur dans les 0.1 et 1 seconde, et dans les 10 secondes~?
\item Quelle part des ordres de la journée simulée est fugace (annulée dans les deux secondes)~?
\end{enumerate}

\medskip\noindent\textbf{Partie IV --- Le verdict.}
\begin{enumerate}[resume]
\item En combien de temps le carnet simulé se reconstitue-t-il après une transaction qui vide son meilleur niveau~?
\item Comment les annulations par transaction de la journée simulée se comparent-elles à MIDAS~?
\item Énoncer le \emph{résultat nommé}~: le tableau des faits stylisés, simulés face aux faits publiés, et ceux que le simulateur ne reproduit pas.
\item Quels écarts comptent pour une étude de la position dans la file et des taux d'exécution, et dans quel sens biaisent-ils l'étude~?
\item Quel écart compte pour une étude des grands ordres~?
\item Qu'est-ce qui a changé dans le flux d'ordres américain de 2012 à 2026, d'après les chiffres de MIDAS~?
\item Que faudrait-il modifier en premier dans le simulateur~?
\item En une phrase~: à quoi sert un fait stylisé~?
\end{enumerate}
\end{problem}

\section{Questions d'entretien}

\begin{interviewq}[$\star$ \iqroles{chercheur}]\label{iq:mx:empirical-facts-of-order-books:1}
Pourquoi le carnet moyen est-il plus profond un ou deux ticks derrière le meilleur prix qu'au meilleur prix~?
\end{interviewq}

\begin{interviewq}[$\star\star$ \iqroles{chercheur, trader}]\label{iq:mx:empirical-facts-of-order-books:2}
Entre actions, les fourchettes sont fortement corrélées à la \omterm{def:mx:empirical-facts-of-order-books:vol}{volatilité par transaction}. Qu'est-ce que cela vous apprend sur ce qui fixe la fourchette~?
\end{interviewq}

\begin{interviewq}[$\star\star$ \iqroles{chercheur}]\label{iq:mx:empirical-facts-of-order-books:3}
La plupart des ordres sont annulés dans la seconde. Comment estimeriez-vous la distribution des durées de vie des ordres lorsque certains ordres sont d'abord exécutés~?
\end{interviewq}

\begin{interviewq}[$\star\star$ \iqroles{développeur, chercheur}]\label{iq:mx:empirical-facts-of-order-books:4}
On vous remet un simulateur de flux d'ordres. Comment décideriez-vous s'il est assez réaliste pour réaliser un backtest d'une stratégie de market making~?
\end{interviewq}

\begin{interviewq}[$\star\star$ \iqroles{trader}]\label{iq:mx:empirical-facts-of-order-books:5}
Pourquoi les fourchettes sont-elles plus larges à l'ouverture qu'à midi~?
\end{interviewq}

\begin{interviewq}[$\star\star\star$ \iqroles{chercheur}]\label{iq:mx:empirical-facts-of-order-books:6}
Que devrait révéler un \omterm{def:mx:empirical-facts-of-order-books:life}{ordre fugace} sur son émetteur, et comment le testeriez-vous~?
\end{interviewq}
